\chapter{Meging and Maxing}
\label{chap:megmax}

\section{Mega forms \includegraphics[width=1em,height=1em]{images/mega.png}}
\label{section:mega}
\input{out/mega}

\section{Primal forms \includegraphics[width=1em,height=1em]{images/primal.png}}
\label{section:primal}

\section{Dynamax forms \includegraphics[width=1em,height=1em]{images/dynamax.png}}
\label{section:dmax}
Dynamax forms have access to three Max Moves: Max Attack, Max Guard, and Max Spirit.
Max Moves have four levels: Locked, 1, 2, and 3.
A freshly acquired (via capture or trade) Dynamax Pokémon has Max Attack at level 1,
  while Max Guard and Max Spirit are both locked.
The Max Attack known depends on the Pokémon's fast attack.
Max Move levels persist across evolution and use of Fast TMs.
\textbf{FIXME: how/when Max Moves are used}
\input{out/dynas}

\section{Gigantamax forms \includegraphics[width=1em,height=1em]{images/gigantamax.png}}
\label{section:gmax}
Gigantamax forms are similar to Dynamax forms to the degree that they can
  employ Max Moves in Max Battles.
Unlike Dynamax forms, each Gigantamax knows a species-specific attack
  called a G-Max Move, which does not change to match the fast attack.
G-Max attacks of level 1, 2, and 3 have power of 350, 400, and 450
  respectively, 100 more than Dynamax attacks of the same level.
Gigantamax forms have their own visual representation.
\input{out/gigantas}
